```latex
\documentclass[11pt]{article}

\usepackage{amsmath,amssymb,amsthm,enumitem,mathtools}

\newtheorem{theorem}{Theorem}
\newtheorem{lemma}[theorem]{Lemma}
\newtheorem{proposition}[theorem]{Proposition}
\newtheorem{corollary}[theorem]{Corollary}

\begin{document}

\title{A counterexample to derived fully faithfulness for the functor \(\pi^{-1}\)}
\author{}
\date{}
\maketitle

\begin{theorem}
The answer to the question in the problem statement is \emph{no}. More precisely, there exists a monoid \(P\) and a presheaf of abelian groups \(F\) on \(B_P^{\mathrm{op}}\) such that the natural morphism
\[
F \longrightarrow \mathrm R\pi_\ast \pi^{-1}F
\]
is not a quasi-isomorphism. In fact, this already fails in degree \(0\).
\end{theorem}

\begin{proof}
We divide the proof into several steps.

\medskip
\noindent
\textbf{1. The functor \(\pi_\ast\) is exact.}

Let
\[
i \colon B_P^{\mathrm{op}} \hookrightarrow A_P^{\mathrm{op}}
\]
be the inclusion of sites. Since \(B_P^{\mathrm{op}}\) carries the trivial topology, every presheaf on \(B_P^{\mathrm{op}}\) is a sheaf. Thus
\[
\operatorname{Sh}(B_P)=\operatorname{PSh}(B_P^{\mathrm{op}}).
\]
By definition,
\[
\pi^{-1}=a\circ \operatorname{Lan}_i,
\]
where \(a\) is sheafification on \(A_P^{\mathrm{op}}\), and \(\operatorname{Lan}_i\) is left Kan extension on presheaves.

For any \(F\in \operatorname{PSh}(B_P^{\mathrm{op}})\) and any sheaf \(G\) on \(A_P^{\mathrm{op}}\),
\[
\operatorname{Hom}_{\operatorname{Sh}(A_P)}(\pi^{-1}F,G)
=
\operatorname{Hom}_{\operatorname{PSh}(A_P^{\mathrm{op}})}(\operatorname{Lan}_iF,G)
=
\operatorname{Hom}_{\operatorname{PSh}(B_P^{\mathrm{op}})}(F,i^\ast G).
\]
Hence the right adjoint \(\pi_\ast\) is simply restriction:
\[
\pi_\ast = i^\ast.
\]
Restriction is exact (it is computed objectwise), so
\[
\mathrm R\pi_\ast = \pi_\ast.
\]
Therefore it is enough to show that, for a suitable \(P\) and \(F\), the ordinary unit
\[
F \longrightarrow \pi_\ast \pi^{-1}F
\]
is not an isomorphism.

\medskip
\noindent
\textbf{2. The monoid \(P=\mathbf N^2\) and a birhizomic cover of \(P\).}

Set
\[
P=\mathbf N e_1 \oplus \mathbf N e_2 \cong \mathbf N^2.
\]
We write monoids additively.

Define submonoids \(Q_1,Q_2\subset \mathbf Z^2\) by
\[
Q_1 := \mathbf N\langle e_1,e_2,\delta_1\rangle
\qquad\text{with }\delta_1:=(-1,2),
\]
\[
Q_2 := \mathbf N\langle e_1,e_2,\delta_2\rangle
\qquad\text{with }\delta_2:=(2,-1).
\]
Equivalently,
\[
Q_1 \cong (P\oplus \mathbf N\delta_1)/(\delta_1+e_1=2e_2),
\qquad
Q_2 \cong (P\oplus \mathbf N\delta_2)/(\delta_2+e_2=2e_1).
\]

\begin{lemma}\label{lemma:Qisfs}
The monoids \(Q_1\) and \(Q_2\) are integral, saturated, and sharp. Moreover the maps
\[
f_1\colon P\to Q_1,\qquad f_2\colon P\to Q_2
\]
are of finite presentation.
\end{lemma}

\begin{proof}
We treat \(Q_1\); the proof for \(Q_2\) is symmetric.

Let \(C_1\subset \mathbf R^2\) be the cone generated by \(e_1=(1,0)\) and \(\delta_1=(-1,2)\). Then
\[
C_1=\{(m,n)\in \mathbf R^2 : n\ge 0,\; 2m+n\ge 0\}.
\]
We claim that
\[
Q_1 = C_1\cap \mathbf Z^2.
\]
Indeed, the inclusion \(Q_1\subseteq C_1\cap \mathbf Z^2\) is clear.

Conversely, let \((m,n)\in C_1\cap \mathbf Z^2\). Then \(n\ge 0\) and \(2m+n\ge 0\).

\begin{itemize}
\item If \(n\) is even, say \(n=2r\), then
\[
(m,n)=(m+r)e_1+r\delta_1,
\]
and \(m+r=\frac{2m+n}{2}\ge 0\), so \((m,n)\in Q_1\).

\item If \(n\) is odd, say \(n=2r+1\), then
\[
(m,n)=(m+r)e_1+r\delta_1+e_2,
\]
and again \(m+r=\frac{2m+n-1}{2}\ge 0\), so \((m,n)\in Q_1\).
\end{itemize}

Thus \(Q_1=C_1\cap \mathbf Z^2\). Since it is the monoid of lattice points in a rational polyhedral cone, it is integral and saturated; and since \(C_1\) is strongly convex, \(Q_1\) is sharp. The displayed presentation shows that \(f_1\) is finitely presented. The same argument applies to \(Q_2\).
\end{proof}

\begin{proposition}\label{prop:cover}
The family
\[
\{\,f_1\colon P\to Q_1,\; f_2\colon P\to Q_2\,\}
\]
is a birhizomic cover of \(P\).
\end{proposition}

\begin{proof}
Let \(V\) be a valuative monoid and let
\[
v\colon P\to V
\]
be a monoid homomorphism. Put
\[
a:=v(e_1),\qquad b:=v(e_2)\in V.
\]
Since \(V\) is valuative, either \(b\ge a\) or \(a\ge b\).

\smallskip
\noindent
\emph{Case 1: \(b\ge a\).}
Then \(b-a\in V\), hence
\[
2b-a=b+(b-a)\in V.
\]
So there is a unique monoid homomorphism
\[
g_1\colon Q_1\to V
\]
with \(g_1\circ f_1=v\), namely the one sending \(\delta_1\) to \(2b-a\).

We claim that \(f_1\) is birhizomic at \(g_1\). By Lemma~\ref{lemma:Qisfs}, \(f_1\) is of finite presentation, so it remains to verify the lifting property.

Consider a factorization
\[
P \xrightarrow{h} M \xrightarrow{p} V
\]
of \(v\), where \(p\) is infinitesimal. Set
\[
m_1:=h(e_1),\qquad m_2:=h(e_2)\in M,
\]
and define
\[
z:=2m_2-m_1\in M^{\mathrm{gp}},\qquad y:=m_1+m_2\in M.
\]
Then
\[
z+y=3m_2\in M,\qquad 2y-z=3m_1\in M,
\]
so
\[
-y\le z\le 2y.
\]
Hence \(z\prec y\).

Moreover
\[
p(z)=2b-a\in V.
\]
Also
\[
2p(z)-p(y)=2(2b-a)-(a+b)=3(b-a)\in V,
\]
so
\[
0\cdot p(z)\le p(y)\le 2p(z).
\]
Thus
\[
p(y)\prec p(z).
\]
Since \(p\) is infinitesimal, from
\[
p(z)\in V,\qquad z\prec y,\qquad p(y)\prec p(z)
\]
we conclude that \(z\in M\).

Now the defining relation \(\delta_1+e_1=2e_2\) shows that there is a unique monoid homomorphism
\[
u\colon Q_1\to M
\]
sending \(\delta_1\mapsto z\) and extending \(h\). By construction,
\[
u\circ f_1=h,\qquad p\circ u=g_1.
\]
Uniqueness is clear because \(Q_1\) is generated by the image of \(P\) and \(\delta_1\).

So \(f_1\) is birhizomic at \(g_1\).

\smallskip
\noindent
\emph{Case 2: \(a\ge b\).}
By symmetry there is a unique \(g_2\colon Q_2\to V\) with \(g_2\circ f_2=v\), and the same argument (interchanging \(e_1\) and \(e_2\)) shows that \(f_2\) is birhizomic at \(g_2\).

\smallskip
Thus every valuative map \(v\colon P\to V\) factors through one of \(Q_1\) or \(Q_2\), with the required birhizomicity. Therefore \(\{f_1,f_2\}\) is a birhizomic cover of \(P\).
\end{proof}

\medskip
\noindent
\textbf{3. The induced cover on \(B_P^{\mathrm{op}}\).}

Let
\[
X:=P,\qquad D_1:=P[-e_1]^\sharp,\qquad D_2:=P[-e_2]^\sharp,\qquad \eta:=P[-(e_1+e_2)]^\sharp.
\]
For \(P=\mathbf N^2\), every object of \(B_P\) is isomorphic to one of \(X,D_1,D_2,\eta\): indeed, for \(x=(m,n)\in \mathbf N^2\),
\[
P[-x]^\sharp \cong
\begin{cases}
X,& x=(0,0),\\
D_1,& m>0,\ n=0,\\
D_2,& m=0,\ n>0,\\
\eta,& m>0,\ n>0.
\end{cases}
\]
Moreover \(\eta\) is the zero monoid.

Let \(J_{\mathrm{ind}}\) denote the Grothendieck topology on \(B_P^{\mathrm{op}}\) induced from \(A_P^{\mathrm{op}}\) by the inclusion \(i\). By definition, the restriction of every sheaf on \(A_P^{\mathrm{op}}\) to \(B_P^{\mathrm{op}}\) is a \(J_{\mathrm{ind}}\)-sheaf.

\begin{lemma}\label{lemma:pullback-cover}
The birhizomic cover \(\{Q_1\to X,\;Q_2\to X\}\) on \(A_P^{\mathrm{op}}\) induces the covering family
\[
\{\,D_1\to X,\; D_2\to X\,\}
\]
on \(B_P^{\mathrm{op}}\).
\end{lemma}

\begin{proof}
In the category \(A_P\), there is a morphism \(Q_1\to D_1\): writing \(D_1\cong \mathbf N\) with generator \(\overline e_2\), define
\[
e_1\mapsto 0,\qquad e_2\mapsto \overline e_2,\qquad \delta_1\mapsto 2\overline e_2.
\]
This respects the relation \(\delta_1+e_1=2e_2\). Similarly, there is a morphism \(Q_2\to D_2\).

There is no morphism \(Q_1\to X\): such a morphism would have to send \(\delta_1\) to an element \(u\in P\) with
\[
u+e_1=2e_2,
\]
i.e.
\[
u=2e_2-e_1\notin P,
\]
impossible. Similarly, there is no morphism \(Q_2\to X\).

There is no morphism \(Q_1\to D_2\): in \(D_2\), the image of \(e_1\) is nonzero and the image of \(e_2\) is \(0\), so the relation
\[
\delta_1+e_1=2e_2
\]
would force
\[
\varphi(\delta_1)+\varphi(e_1)=0.
\]
But in a sharp monoid, if \(x+y=0\), then \(x=y=0\); thus \(\varphi(e_1)=0\), contradiction. Likewise there is no morphism \(Q_2\to D_1\).

Finally, both \(Q_1\) and \(Q_2\) map to \(\eta\) (since \(\eta\) is the zero monoid), and \(\eta\) already lies in the sieve generated by \(D_1\to X\) and \(D_2\to X\).

Thus the inverse-image sieve on \(B_P^{\mathrm{op}}\) of the covering sieve generated by \(Q_1\to X\) and \(Q_2\to X\) is exactly the sieve generated by \(D_1\to X\) and \(D_2\to X\). Hence \(\{D_1\to X,\;D_2\to X\}\) is a \(J_{\mathrm{ind}}\)-cover.
\end{proof}

\begin{lemma}\label{lemma:fiberproduct}
In \(B_P^{\mathrm{op}}\), the fiber product of \(D_1\to X\) and \(D_2\to X\) is \(\eta\).
\end{lemma}

\begin{proof}
Fiber products in \(B_P^{\mathrm{op}}\) are pushouts in \(B_P\). The pushout of
\[
P\to P[-e_1]^\sharp,\qquad P\to P[-e_2]^\sharp
\]
is obtained by forcing both \(e_1\) and \(e_2\) to become \(0\); this is exactly
\[
P[-(e_1+e_2)]^\sharp=\eta.
\]
\end{proof}

\medskip
\noindent
\textbf{4. A presheaf on \(B_P^{\mathrm{op}}\) which is not a \(J_{\mathrm{ind}}\)-sheaf.}

Define a presheaf \(F\) of abelian groups on \(B_P^{\mathrm{op}}\) by
\[
F(P[-x]^\sharp)=
\begin{cases}
\mathbf Z,& x=0,\\
0,& x\neq 0.
\end{cases}
\]
All nonidentity transition maps are necessarily \(0\).

Thus
\[
F(X)=\mathbf Z,\qquad F(D_1)=F(D_2)=F(\eta)=0.
\]

By Lemma~\ref{lemma:pullback-cover}, \(\{D_1\to X,\;D_2\to X\}\) is a \(J_{\mathrm{ind}}\)-cover, and by Lemma~\ref{lemma:fiberproduct} its overlap is \(\eta\). Therefore a \(J_{\mathrm{ind}}\)-sheaf \(G\) must satisfy
\[
G(X)\xrightarrow{\sim} G(D_1)\times_{G(\eta)}G(D_2).
\]
For our \(F\), the right-hand side is
\[
F(D_1)\times_{F(\eta)}F(D_2)=0\times_0 0=0,
\]
while \(F(X)=\mathbf Z\). Hence \(F\) is \emph{not} a \(J_{\mathrm{ind}}\)-sheaf.

On the other hand, \(G:=\pi_\ast\pi^{-1}F\) is the restriction to \(B_P^{\mathrm{op}}\) of the sheaf \(\pi^{-1}F\) on \(A_P^{\mathrm{op}}\), so \(G\) is a \(J_{\mathrm{ind}}\)-sheaf by construction of the induced topology.

Therefore \(F\not\cong G\) as presheaves on \(B_P^{\mathrm{op}}\); equivalently,
\[
F \longrightarrow \pi_\ast\pi^{-1}F
\]
is not an isomorphism.

Since \(\mathrm R\pi_\ast=\pi_\ast\), the map
\[
F\longrightarrow \mathrm R\pi_\ast\pi^{-1}F
\]
is not a quasi-isomorphism.

This proves the theorem.
\end{proof}

\begin{corollary}
Already for \(P=\mathbf N^2\), the topology induced on \(B_P^{\mathrm{op}}\) from the birhizomic topology on \(A_P^{\mathrm{op}}\) is strictly finer than the trivial topology.
\end{corollary}

\end{document}
```